% ------------------------------------------------------------------------
% Configuração e Arquitetura do Template LaTeX
% Desenvolvido por: Aurora Drumond Costa Magalhães
% Ano: 2026
% ------------------------------------------------------------------------

\documentclass[
    12pt,
    openright,
    oneside,
    a4paper,
    english,
    brazil
    ]{abntex2}

% --- PACOTES GERAIS ---
\usepackage{lmodern}
\usepackage[T1]{fontenc}
\usepackage[utf8]{inputenc}
\usepackage{indentfirst}
\usepackage{graphicx}
\usepackage{float}
\usepackage[alf]{abntex2cite}

% --- CONFIGURAÇÃO DO AMBIENTE QUADRO ---
\newcommand{\quadroname}{Quadro}
\newcommand{\listofquadrosname}{Lista de quadros}

% Utilizando a sintaxe estrita do pacote 'float' (\newfloat{nome}{posicionamento}{extensão}[contador])
\floatstyle{plaintop} % Garante que a legenda fique acima do quadro
\newfloat{quadro}{htb}{loq}[chapter]
\floatname{quadro}{\quadroname}

% Cria a lista de quadros associada ao sumário
\newlistof{listofquadros}{loq}{\listofquadrosname}
\newlistentry{quadro}{loq}{0}

% Define a numeração contínua (evita Quadro 1.1, Quadro 1.2)
\counterwithout{quadro}{chapter}

% Configura a formatação de exibição da Lista de Quadros no PDF
\renewcommand{\cftquadroname}{\quadroname\space}
\renewcommand*{\cftquadroaftersnum}{\hfill--\hfill}

% --- IMPORTA SEUS ESTILOS PERSONALIZADOS ---
% Substitua pelos caminhos reais dos seus arquivos de estilo (.sty), caso possua.
\usepackage{styles/code-style} 
\usepackage{styles/mylib-style}

% --- DADOS DA CAPA E FOLHA DE ROSTO ---
\titulo{Poke-Adocao: Sistema de Adoção de Pokémon via Geolocalização}
\autor{Aurora Drumond Costa Magalhães \and Ana Clara de Souza \and Kayke Wellington}
\local{Timóteo - MG}
\data{2026}

\instituicao{%
  Centro Universitário Católica do Leste de Minas Gerais
  \par
  Coronel Fabriciano}

\tipotrabalho{Documentação de Arquitetura de Software}

\preambulo{Relatório técnico apresentado como documentação oficial da arquitetura do sistema Poke-Adocao, detalhando as especificações de engenharia do serviço de backend e da aplicação móvel.}

% --- INÍCIO DO DOCUMENTO ---
\begin{document}

% Elementos Pré-Textuais (Obrigatórios segundo a ABNT)
\imprimircapa
\imprimirfolhaderosto

% --- RESUMOS ---
\setlength{\absparsep}{18pt}

\begin{resumo}
O projeto Poke-Adocao consiste num sistema de adoção de Pokémon que utiliza geolocalização para interligar os utilizadores a criaturas virtuais no espaço físico. A arquitetura baseia-se num monorepo que integra uma API RESTful em Python com FastAPI no backend, e uma aplicação móvel multiplataforma construída em React Native e Expo no frontend. O núcleo do processo de adoção é gerido por uma Máquina de Estados Finitos (FSM). Atualmente operando de forma binária para garantir a integridade do ciclo de vida restrito a um limite rigoroso de 50 metros de distância, a FSM demonstra a fundação arquitetural necessária para suportar, futuramente, estados intermediários mais dinâmicos e contratos \textit{player-to-player} (P2P) complexos, como transações pendentes de aprovação. Para viabilizar este mecanismo, o sistema implementa um motor espacial que calcula distâncias geográficas utilizando a fórmula de Haversine diretamente na camada da aplicação, processo previamente otimizado por uma filtragem através de \textit{Bounding Box}. A persistência de dados é assegurada de forma leve e contida por um banco de dados relacional SQLite, escolhido estrategicamente para validar as transações de \textit{optimistic locking} e o fluxo de ponta a ponta do protótipo sem a sobrecarga de configuração de infraestrutura externa. O presente documento detalha as decisões de engenharia, a infraestrutura e os requisitos técnicos que sustentam a aplicação.

\vspace{\onelineskip}
\noindent
\textbf{Palavras-chave:} Geolocalização. React Native. FastAPI. Máquina de Estados Finitos. Haversine. SQLite.
\end{resumo}

\begin{resumo}[Abstract]
\begin{otherlanguage*}{english}
The Poke-Adocao project is a Pokemon adoption system that leverages geolocation to connect users with virtual creatures in the physical world. The architecture is built upon a monorepo integrating a RESTful API developed in Python with FastAPI for the backend, and a cross-platform mobile application built with React Native and Expo for the frontend. The core adoption process is managed by a Finite State Machine (FSM). Currently operating in a binary manner to enforce the lifecycle under a strict 50-meter geographic threshold, the FSM demonstrates the architectural foundation necessary to support complex intermediate states and dynamic player-to-player (P2P) contracts in the future, such as pending approvals. To enable this mechanism, the system implements a spatial engine that computes geographic distances using the Haversine formula at the application layer, optimized by a Bounding Box pre-filtering process. Data persistence is handled in a lightweight and contained manner by a relational SQLite database, strategically chosen to validate optimistic locking transactions and the end-to-end flow of the prototype without the overhead of external infrastructure configuration. This document details the engineering decisions, infrastructure, and technical requirements underlying the application.

\vspace{\onelineskip}
\noindent
\textbf{Keywords:} Geolocation. React Native. FastAPI. Finite State Machine. Haversine. SQLite.
\end{otherlanguage*}
\end{resumo}

% --- LISTAS ---
% Descomente as listas (removendo o sinal '%') que forem necessárias para o seu trabalho de acordo com os elementos presentes.

% Lista de Figuras
\pdfbookmark[0]{\listfigurename}{lof}
\listoffigures*
\cleardoublepage

% Lista de Códigos / Trechos de Software
% \pdfbookmark[0]{\lstlistlistingname}{lol}
% \lstlistoflistings
% \cleardoublepage

% Lista de Quadros
\pdfbookmark[0]{Lista de quadros}{loq}
\listofquadros*
\cleardoublepage

% Sumário (Obrigatório)
\pdfbookmark[0]{\contentsname}{toc}
\tableofcontents
\cleardoublepage

% --------------------------------------------------------
% ELEMENTOS TEXTUAIS
% --------------------------------------------------------
\textual

% --- ESTRUTURAÇÃO EM MÚLTIPLOS ARQUIVOS ---
% A chamada \include insere o conteúdo de arquivos externos. Isso permite modularizar o 
% documento separando os tópicos exigidos no relatório técnico em arquivos .tex 
% distintos dentro do diretório 'capitulos/'.
% Não é necessário adicionar a extensão .tex ao usar o comando \include.

\chapter{Design e Arquitetura}

O padrão arquitetural macro adotado pelo sistema Poke-Adocao fundamenta-se no modelo Cliente-Servidor Desacoplado, operacionalizado sob a estrutura de um monorrepositório. Essa abordagem centraliza a governança do código-fonte e dos contratos de comunicação, permitindo o isolamento completo entre a aplicação cliente móvel e a API de serviços, que interagem estritamente por meio de protocolos de rede \textit{stateless}.

\section{Visão Lógica e de Componentes}

A organização interna do ecossistema é bifurcada entre as camadas funcionais de \textit{backend} e \textit{frontend}, estruturadas de forma modular para assegurar a separação de conceitos \cite{sommerville2011software}. 

No \textit{backend}, a API segue uma Arquitetura em Camadas dirigida por serviços. A camada de roteamento e entrada, situada no arquivo \texttt{main.py}, atua como a interface de exposição dos \textit{endpoints} HTTP, gerenciando \textit{middlewares} e a orquestração de controladores assíncronos. A validação dos dados é centralizada no \texttt{schemas.py}, que utiliza o \textit{Pydantic} para garantir a tipagem estática e serialização de todas as requisições. As regras de negócio operam em uma camada dedicada composta por serviços como \texttt{adoption\_service.py} e \texttt{spatial\_service.py}, que isolam a lógica da Máquina de Estados Finitos (FSM) e as validações trigonométricas, enquanto a persistência é orquestrada através dos modelos declarativos em \texttt{models.py} e \texttt{database.py}, gerindo o ciclo de vida das sessões transacionais.

No \textit{frontend}, a interface adota uma arquitetura baseada em componentes funcionais alimentados por injeção de estado via \textit{hooks}. O \texttt{RootNavigator.tsx} gerencia o fluxo global de autenticação condicional, chaveando o utilizador entre a tela de \textit{login} e a navegação principal estruturada no \texttt{MainTabs.tsx}. A gestão de estado é centralizada nativamente no \texttt{AuthContext.tsx}, encapsulando a persistência de \textit{tokens} e simplificando a telemetria, que é isolada no \textit{hook} reativo \texttt{useUserLocation.ts}, abstraindo a complexidade dos sensores de geolocalização nativos.

\section{Visão de Dados}

O modelo de persistência baseia-se no paradigma Relacional para salvaguardar a consistência da FSM. Atualmente, o protótipo utiliza o \textit{SGBD SQLite} para validar o fluxo de ponta a ponta com \textit{optimistic locking}, mantendo compatibilidade arquitetural para migração para \textit{PostgreSQL} em escala de produção. O mapeamento é orquestrado via \textit{SQLAlchemy}, estruturando o esquema lógico em torno de três entidades fundamentais: \texttt{users}, para metadados de identificação; \texttt{pokemons}, para referências da PokeAPI e estado da entidade; e \texttt{adoptions}, uma tabela de junção que materializa o vínculo transacional e audita o ciclo de vida da adoção.

\section{Visão Física e de Implantação}

A infraestrutura distribui as cargas de processamento em três nós lógicos de implantação. A Camada Cliente consiste nos artefatos compilados via \textit{Expo}, que executam nativamente em dispositivos móveis consumindo o hardware do receptor GPS. A Camada de           
\chapter{Explicação de Trechos Importantes de Código}

A implementação do sistema fundamenta-se na integração entre a captura de sensores nativos no \textit{frontend} e a validação estrita de regras de negócio no \textit{backend}. Abaixo, detalham-se os componentes críticos que operacionalizam a telemetria, o gerenciamento de identidade e a validação espacial.

\section{Captura e Telemetria de Geolocalização}

A interface móvel abstrai a complexidade do hardware via \textit{hooks} customizados. O \texttt{useUserLocation.ts} implementa a observação passiva de coordenadas, integrando o \texttt{expo-location} para garantir a precisão necessária às validações da FSM. A escolha pela precisão \texttt{High} e um intervalo de \texttt{timeInterval: 5000} equilibra a responsividade do radar com o consumo de energia do dispositivo.

\begin{lstlisting}[style=estiloJS, caption={Hook de telemetria para captura de geolocalização}]
export const useUserLocation = () => {
  const [location, setLocation] = useState<Location.LocationObject | null>(null);

  useEffect(() => {
    (async () => {
      let { status } = await Location.requestForegroundPermissionsAsync();
      if (status !== 'granted') return;

      await Location.watchPositionAsync({
        accuracy: Location.Accuracy.High,
        timeInterval: 5000,
      }, (newLocation) => {
        setLocation(newLocation);
      });
    })();
  }, []);

  return location;
};
\end{lstlisting}

\section{Gerenciamento de Autenticação e Navegação}

O controle de sessão é consolidado no \texttt{AuthContext.tsx}, que atua como o \textit{Provider} central. A arquitetura evita a propagação excessiva de propriedades (\textit{prop-drilling}) ao injetar o estado do usuário diretamente na árvore de navegação. No \texttt{RootNavigator.tsx}, a condicional de renderização assegura que o estado de autenticação dite a composição da \textit{stack} de navegação, roteando o usuário dinamicamente entre as telas de \textit{login} e a aplicação funcional (\texttt{MainTabs.tsx}).

\begin{lstlisting}[style=estiloJS, caption={Navegação condicional baseada no contexto de autenticação}]
const RootNavigator = () => {
  const { user } = useContext(AuthContext);

  return (
    <NavigationContainer>
      {user ? <MainTabs /> : <LoginScreen />}
    </NavigationContainer>
  );
};
\end{lstlisting}

\section{Validação Espacial no Servidor}

Diferente do \textit{frontend}, que provê a telemetria, a inteligência do sistema reside no \texttt{spatial\_service.py}. A validação de proximidade para a adoção de um Pokémon não confia no dispositivo do usuário; o servidor recalcula a distância geodésica em relação à coordenada absoluta da entidade persistida no banco.

\begin{lstlisting}[style=estiloPython, caption={Validação de proximidade via Haversine no servidor}]
def is_within_distance(lat1, lon1, lat2, lon2, limit=50):
    dlat = radians(lat2 - lat1)
    dlon = radians(lon2 - lon1)
    a = sin(dlat/2)**2 + cos(radians(lat1)) * cos(radians(lat2)) * sin(dlon/2)**2
    c = 2 * atan2(sqrt(a), sqrt(1-a))
    distance = 6371000 * c  # Distancia em metros
    return distance <= limit
\end{lstlisting}

Este trecho é o núcleo da segurança da FSM. Ao executar esta validação isoladamente no servidor, o sistema blinda-se contra \textit{spoofing} de GPS, garantindo que o estado de \textit{ADOPTED} só seja atingido quando a proximidade física for matematicamente comprovada pela fórmula de Haversine.            
\chapter{Testes e Validação}

Este capítulo detalha as estratégias de garantia de qualidade (QA) e as metodologias de validação aplicadas para atestar a integridade, o desempenho e a segurança do ecossistema Poke-Adocao. O processo de homologação afasta-se de validações empíricas manuais em prol de uma esteira de testes automatizados com cobertura arquitetural contínua.

\section{Cenários de Teste e Automação}

A validação do fluxo transacional de ponta a ponta (\textit{End-to-End}) é assegurada por testes unitários e de integração rigorosos, divididos entre os dois pólos do monorrepositório, isolando as responsabilidades de estado e de interface.

\subsection{Validação do \textit{Backend}}
No ambiente de servidor, a infraestrutura de testes é orquestrada pela \textit{framework} \textit{Pytest}, acoplada à extensão \textit{pytest-asyncio} para suportar a resolução do \textit{event loop} do FastAPI. O banco de dados relacional é interceptado, e as transações de teste ocorrem em bases \textit{SQLite} em memória para evitar poluição de dados persistentes.
Os cenários homologados na base de código garantem as invariantes do domínio espacial e transacional:
\begin{itemize}
    \item \textbf{Motor Espacial (\texttt{test\_spatial\_service.py}):} Testes focados em atestar a precisão matemática da Fómula de Haversine e a delimitação dos vértices da \textit{Bounding Box} contra anomalias trigonométricas de latitude e longitude.
    \item \textbf{Máquina de Estados (\texttt{test\_adoption\_service.py}):} Validação da transição unidirecional da \textit{AdoptionStatus}, forçando falhas deliberadas como adoção concorrente da mesma entidade (\textit{race conditions}), violação da distância máxima de 50 metros e excedente do limite de inventário da \textit{party}.
    \item \textbf{Contratos de API (\texttt{test\_api.py} e \texttt{test\_auth.py}):} Verificação da emissão de \textit{tokens} criptográficos, proteção de rotas restritas e sanitização de requisições HTTP malformadas.
\end{itemize}

\subsection{Validação do \textit{Frontend}}
Na camada cliente, a integridade da aplicação móvel \textit{React Native} é aferida através do motor \textit{Jest} e das bibliotecas de dom \textit{React Native Testing Library}.
Os testes isolam o comportamento da interface do usuário da latência de rede:
\begin{itemize}
    \item \textbf{Comportamento de Hooks (\texttt{useUserLocation.test.ts}):} Simulação do ciclo de vida dos sensores do dispositivo móvel para assegurar que a telemetria geográfica é corretamente capturada, filtrada e exposta ao contexto da aplicação.
    \item \textbf{Renderização de Telas (\texttt{LoginScreen}, \texttt{PartyScreen}, \texttt{ProfileScreen}):} Atesta-se a montagem da árvore de componentes, garantindo que o inventário (\textit{Party}) não renderiza instâncias nulas e que o roteamento restringe utilizadores não autenticados.
\end{itemize}

\section{Visão do Usuário}

\begin{figure}[htb]
    \centering
    \caption{Tela de Login}
    \label{fig:login}
    \includegraphics[width=\linewidth, height=0.6\textheight, keepaspectratio]{figure/login.jpeg} 
    \fonte{Elaborado pelos autores (2026).}
\end{figure}

\begin{figure}[htb]
    \centering
    \caption{Tela do grupo de pokemons do usuário (party)}
    \label{fig:party}
    \includegraphics[width=\linewidth, height=0.6\textheight, keepaspectratio]{figure/party.jpeg} 
    \fonte{Elaborado pelos autores (2026).}
\end{figure}

\begin{figure}[htb]
    \centering
    \caption{Tela da lista de adoções vazia}
    \label{fig:lista_de_adocao}
    \includegraphics[width=\linewidth, height=0.6\textheight, keepaspectratio]{figure/lista_de_adocao.jpeg} 
    \fonte{Elaborado pelos autores (2026).}
\end{figure}

\begin{figure}[htb]
    \centering
    \caption{Tela da lista de adoções populada}
    \label{fig:lista_de_adocao_populada}
    \includegraphics[width=\linewidth, height=0.6\textheight, keepaspectratio]{figure/populada.jpeg} 
    \fonte{Elaborado pelos autores (2026).}
\end{figure}

\begin{figure}[htb]
    \centering
    \caption{Tela do Radar (Sem Pokemons ou jogadores perto)}
    \label{fig:radar}
    \includegraphics[width=\linewidth, height=0.6\textheight, keepaspectratio]{figure/radar (indisponivel).jpeg} 
    \fonte{Elaborado pelos autores (2026).}
\end{figure}

\begin{figure}[htb]
    \centering
    \caption{Tela do treinador (parte 1)}
    \label{fig:treinadorum}
    \includegraphics[width=\linewidth, height=0.6\textheight, keepaspectratio]{figure/treinador.jpeg} 
    \fonte{Elaborado pelos autores (2026).}
\end{figure}

\begin{figure}[htb]
    \centering
    \caption{Tela do treinador (parte 2)}
    \label{fig:treinadordois}
    \includegraphics[width=\linewidth, height=0.6\textheight, keepaspectratio]{figure/treinador (final).jpeg} 
    \fonte{Elaborado pelos autores (2026).}
\end{figure}           
\chapter{Conclusão}

O desenvolvimento do sistema Poke-Adocao permitiu a concretização da arquitetura de geolocalização e gerenciamento de estado proposta inicialmente. Os objetivos gerais e específicos delineados no Capítulo 3 foram integralmente alcançados, com a entrega de um sistema funcional que integra uma API robusta de \textit{backend} em Python e FastAPI a uma interface móvel multiplataforma em React Native. A validação das regras de negócio, assegurada pela Máquina de Estados Finitos (FSM) e pelo motor espacial baseado na fórmula de Haversine, demonstrou que a arquitetura proposta é capaz de garantir a integridade transacional mesmo sob condições de concorrência e restrições de proximidade física.

Do ponto de vista da engenharia de \textit{software}, os maiores gargalos técnicos enfrentados concentraram-se na mitigação da latência de processamento espacial e na garantia da consistência de dados em ambiente distribuído. A complexidade de calcular distâncias geodésicas em tempo real, aliada à necessidade de evitar condições de corrida (\textit{race conditions}) na adoção de entidades únicas, exigiu a implementação de uma camada de bloqueio otimista (\textit{optimistic locking}) no banco de dados. A transição para um modelo de monorrepositório foi determinante para superar a fricção de desalinhamento de contratos entre as camadas, permitindo que a equipe focasse no refinamento dos algoritmos de validação em vez de na gestão de dependências externas.

Quanto aos trabalhos futuros, o projeto mantém uma lista estratégica de dívidas técnicas e melhorias planejadas para a próxima fase de maturação. O plano de expansão inclui:
\begin{itemize}
    \item \textbf{Escalabilidade Geográfica:} Implementação de índices espaciais nativos (\textit{PostGIS}) para substituir a pré-filtragem por \textit{Bounding Box} manual, elevando a performance das consultas para cenários de alta densidade de usuários.
    \item \textbf{Dívida Técnica de \textit{Frontend}:} Migração dos componentes de interface para um sistema de temas totalmente tokenizado, aumentando a manutenibilidade do design.
    \item \textbf{Novas Funcionalidades:} Desenvolvimento de um motor de \textit{caching} distribuído (Redis) para persistir temporariamente as telemetrias de localização dos usuários e reduzir o custo de escrita direta no banco relacional.
\end{itemize}

Em suma, o projeto Poke-Adocao demonstra a viabilidade de construir aplicações de \textit{Location-Based Services} complexas utilizando um \textit{stack} tecnológico moderno e rigoroso, estabelecendo uma fundação sólida para futuras iterações e melhorias de escalabilidade de engenharia.   
\include{capitulos/05-trechos-codigo}       
@book{fowler2002patterns,
  title={Patterns of Enterprise Application Architecture},
  author={Fowler, Martin},
  year={2002},
  publisher={Addison-Wesley Longman},
  address={Boston}
}

@article{sinnott1984virtues,
  title={Virtues of the Haversine},
  author={Sinnott, Roger W},
  journal={Sky and Telescope},
  volume={68},
  number={2},
  pages={159},
  year={1984}
}

@phdthesis{fielding2000architectural,
  title={Architectural styles and the design of network-based software architectures},
  author={Fielding, Roy Thomas},
  year={2000},
  school={University of California, Irvine}
}

@book{gurnani2020developing,
  title={Developing Applications with React Native},
  author={Gurnani, Akshat},
  year={2020},
  publisher={Packt Publishing Ltd},
  address={Birmingham}
}

@book{sommerville2011software,
  title={Software Engineering},
  author={Sommerville, Ian},
  year={2011},
  publisher={Addison-Wesley},
  address={Boston}
}

% ==============================================================================
% MANUAL DE OPERAÇÃO: GERENCIAMENTO DE REFERÊNCIAS (BIBTEX + ABNTEX2)
% ==============================================================================
% 
% 1. O CONCEITO DO ARQUIVO .BIB
% No LaTeX, você não digita a lista de referências manualmente no final do texto.
% O arquivo 'referencias.bib' atua como um banco de dados relacional. Você
% cadastra os metadados das fontes nele e faz "consultas" no texto principal.
% Regra de Ouro: Uma fonte cadastrada no .bib SÓ aparecerá no PDF final se
% for explicitamente citada em algum lugar do texto.
%
% 2. COMO CADASTRAR UMA FONTE (No arquivo referencias.bib)
% Cada entrada possui um "tipo" (@book, @article, @misc) e uma "chave" única
% que você inventa para chamá-la depois.
%
% Exemplo de Livro Técnico:
% @book{sommerville2011,
%   author    = {Sommerville, Ian},
%   title     = {Engenharia de software},
%   publisher = {Pearson Prentice Hall},
%   year      = {2011},
%   address   = {São Paulo}
% }
%
% Exemplo de Site, Documentação ou Repositório:
% @misc{fastapi2026,
%   author        = {Ramírez, Sebastián},
%   title         = {FastAPI Documentation},
%   year          = {2026},
%   url           = {https://fastapi.tiangolo.com/},
%   urlaccessdate = {8 jun. 2026}
% }
% Nota: Autores corporativos (ex: Meta, Google, W3C) devem ser envolvidos 
% em chaves duplas no arquivo .bib para que o LaTeX não inverta os nomes. 
% Ex: author = {{Meta Platforms}}
%
% 3. COMO CHAMAR A FONTE NO TEXTO (Nos arquivos .tex)
% O pacote abntex2cite fornece comandos rigorosos para adequação à NBR 10520.
%
% -> Citação Indireta (Ao final da frase, em maiúsculas):
% Comando: \cite{chave}
% Exemplo: A arquitetura baseada em eventos reduz o acoplamento \cite{sommerville2011}.
% PDF renderizado: A arquitetura baseada em eventos reduz o acoplamento (SOMMERVILLE, 2011).
%
% -> Citação Direta/Narrativa (Inserida no fluxo do texto):
% Comando: \citeonline{chave}
% Exemplo: Segundo \citeonline{sommerville2011}, a engenharia de requisitos é vital.
% PDF renderizado: Segundo Sommerville (2011), a engenharia de requisitos é vital.
%
% -> Inserção Forçada (Sem citação textual):
% Comando: \nocite{chave}
% Uso: Insere a fonte na bibliografia final sem precisar colocar um (AUTOR, Ano) no texto.
% Evite usar isso em trabalhos acadêmicos para burlar a fundamentação.
% ==============================================================================     
% ------------------------------------------------------------------------
% Configuração e Arquitetura do Template LaTeX
% Desenvolvido por: Aurora Drumond Costa Magalhães
% Ano: 2026
% ------------------------------------------------------------------------

\documentclass[
    12pt,
    openright,
    oneside,
    a4paper,
    english,
    brazil
    ]{abntex2}

% --- PACOTES GERAIS ---
\usepackage{lmodern}
\usepackage[T1]{fontenc}
\usepackage[utf8]{inputenc}
\usepackage{indentfirst}
\usepackage{graphicx}
\usepackage{float}
\usepackage[alf]{abntex2cite}

% --- CONFIGURAÇÃO DO AMBIENTE QUADRO ---
\newcommand{\quadroname}{Quadro}
\newcommand{\listofquadrosname}{Lista de quadros}

% Utilizando a sintaxe estrita do pacote 'float' (\newfloat{nome}{posicionamento}{extensão}[contador])
\floatstyle{plaintop} % Garante que a legenda fique acima do quadro
\newfloat{quadro}{htb}{loq}[chapter]
\floatname{quadro}{\quadroname}

% Cria a lista de quadros associada ao sumário
\newlistof{listofquadros}{loq}{\listofquadrosname}
\newlistentry{quadro}{loq}{0}

% Define a numeração contínua (evita Quadro 1.1, Quadro 1.2)
\counterwithout{quadro}{chapter}

% Configura a formatação de exibição da Lista de Quadros no PDF
\renewcommand{\cftquadroname}{\quadroname\space}
\renewcommand*{\cftquadroaftersnum}{\hfill--\hfill}

% --- IMPORTA SEUS ESTILOS PERSONALIZADOS ---
% Substitua pelos caminhos reais dos seus arquivos de estilo (.sty), caso possua.
\usepackage{styles/code-style} 
\usepackage{styles/mylib-style}

% --- DADOS DA CAPA E FOLHA DE ROSTO ---
\titulo{Poke-Adocao: Sistema de Adoção de Pokémon via Geolocalização}
\autor{Aurora Drumond Costa Magalhães \and Ana Clara de Souza \and Kayke Wellington}
\local{Timóteo - MG}
\data{2026}

\instituicao{%
  Centro Universitário Católica do Leste de Minas Gerais
  \par
  Coronel Fabriciano}

\tipotrabalho{Documentação de Arquitetura de Software}

\preambulo{Relatório técnico apresentado como documentação oficial da arquitetura do sistema Poke-Adocao, detalhando as especificações de engenharia do serviço de backend e da aplicação móvel.}

% --- INÍCIO DO DOCUMENTO ---
\begin{document}

% Elementos Pré-Textuais (Obrigatórios segundo a ABNT)
\imprimircapa
\imprimirfolhaderosto

% --- RESUMOS ---
\setlength{\absparsep}{18pt}

\begin{resumo}
O projeto Poke-Adocao consiste num sistema de adoção de Pokémon que utiliza geolocalização para interligar os utilizadores a criaturas virtuais no espaço físico. A arquitetura baseia-se num monorepo que integra uma API RESTful em Python com FastAPI no backend, e uma aplicação móvel multiplataforma construída em React Native e Expo no frontend. O núcleo do processo de adoção é gerido por uma Máquina de Estados Finitos (FSM). Atualmente operando de forma binária para garantir a integridade do ciclo de vida restrito a um limite rigoroso de 50 metros de distância, a FSM demonstra a fundação arquitetural necessária para suportar, futuramente, estados intermediários mais dinâmicos e contratos \textit{player-to-player} (P2P) complexos, como transações pendentes de aprovação. Para viabilizar este mecanismo, o sistema implementa um motor espacial que calcula distâncias geográficas utilizando a fórmula de Haversine diretamente na camada da aplicação, processo previamente otimizado por uma filtragem através de \textit{Bounding Box}. A persistência de dados é assegurada de forma leve e contida por um banco de dados relacional SQLite, escolhido estrategicamente para validar as transações de \textit{optimistic locking} e o fluxo de ponta a ponta do protótipo sem a sobrecarga de configuração de infraestrutura externa. O presente documento detalha as decisões de engenharia, a infraestrutura e os requisitos técnicos que sustentam a aplicação.

\vspace{\onelineskip}
\noindent
\textbf{Palavras-chave:} Geolocalização. React Native. FastAPI. Máquina de Estados Finitos. Haversine. SQLite.
\end{resumo}

\begin{resumo}[Abstract]
\begin{otherlanguage*}{english}
The Poke-Adocao project is a Pokemon adoption system that leverages geolocation to connect users with virtual creatures in the physical world. The architecture is built upon a monorepo integrating a RESTful API developed in Python with FastAPI for the backend, and a cross-platform mobile application built with React Native and Expo for the frontend. The core adoption process is managed by a Finite State Machine (FSM). Currently operating in a binary manner to enforce the lifecycle under a strict 50-meter geographic threshold, the FSM demonstrates the architectural foundation necessary to support complex intermediate states and dynamic player-to-player (P2P) contracts in the future, such as pending approvals. To enable this mechanism, the system implements a spatial engine that computes geographic distances using the Haversine formula at the application layer, optimized by a Bounding Box pre-filtering process. Data persistence is handled in a lightweight and contained manner by a relational SQLite database, strategically chosen to validate optimistic locking transactions and the end-to-end flow of the prototype without the overhead of external infrastructure configuration. This document details the engineering decisions, infrastructure, and technical requirements underlying the application.

\vspace{\onelineskip}
\noindent
\textbf{Keywords:} Geolocation. React Native. FastAPI. Finite State Machine. Haversine. SQLite.
\end{otherlanguage*}
\end{resumo}

% --- LISTAS ---
% Descomente as listas (removendo o sinal '%') que forem necessárias para o seu trabalho de acordo com os elementos presentes.

% Lista de Figuras
\pdfbookmark[0]{\listfigurename}{lof}
\listoffigures*
\cleardoublepage

% Lista de Códigos / Trechos de Software
% \pdfbookmark[0]{\lstlistlistingname}{lol}
% \lstlistoflistings
% \cleardoublepage

% Lista de Quadros
\pdfbookmark[0]{Lista de quadros}{loq}
\listofquadros*
\cleardoublepage

% Sumário (Obrigatório)
\pdfbookmark[0]{\contentsname}{toc}
\tableofcontents
\cleardoublepage

% --------------------------------------------------------
% ELEMENTOS TEXTUAIS
% --------------------------------------------------------
\textual

% --- ESTRUTURAÇÃO EM MÚLTIPLOS ARQUIVOS ---
% A chamada \include insere o conteúdo de arquivos externos. Isso permite modularizar o 
% documento separando os tópicos exigidos no relatório técnico em arquivos .tex 
% distintos dentro do diretório 'capitulos/'.
% Não é necessário adicionar a extensão .tex ao usar o comando \include.

\chapter{Design e Arquitetura}

O padrão arquitetural macro adotado pelo sistema Poke-Adocao fundamenta-se no modelo Cliente-Servidor Desacoplado, operacionalizado sob a estrutura de um monorrepositório. Essa abordagem centraliza a governança do código-fonte e dos contratos de comunicação, permitindo o isolamento completo entre a aplicação cliente móvel e a API de serviços, que interagem estritamente por meio de protocolos de rede \textit{stateless}.

\section{Visão Lógica e de Componentes}

A organização interna do ecossistema é bifurcada entre as camadas funcionais de \textit{backend} e \textit{frontend}, estruturadas de forma modular para assegurar a separação de conceitos \cite{sommerville2011software}. 

No \textit{backend}, a API segue uma Arquitetura em Camadas dirigida por serviços. A camada de roteamento e entrada, situada no arquivo \texttt{main.py}, atua como a interface de exposição dos \textit{endpoints} HTTP, gerenciando \textit{middlewares} e a orquestração de controladores assíncronos. A validação dos dados é centralizada no \texttt{schemas.py}, que utiliza o \textit{Pydantic} para garantir a tipagem estática e serialização de todas as requisições. As regras de negócio operam em uma camada dedicada composta por serviços como \texttt{adoption\_service.py} e \texttt{spatial\_service.py}, que isolam a lógica da Máquina de Estados Finitos (FSM) e as validações trigonométricas, enquanto a persistência é orquestrada através dos modelos declarativos em \texttt{models.py} e \texttt{database.py}, gerindo o ciclo de vida das sessões transacionais.

No \textit{frontend}, a interface adota uma arquitetura baseada em componentes funcionais alimentados por injeção de estado via \textit{hooks}. O \texttt{RootNavigator.tsx} gerencia o fluxo global de autenticação condicional, chaveando o utilizador entre a tela de \textit{login} e a navegação principal estruturada no \texttt{MainTabs.tsx}. A gestão de estado é centralizada nativamente no \texttt{AuthContext.tsx}, encapsulando a persistência de \textit{tokens} e simplificando a telemetria, que é isolada no \textit{hook} reativo \texttt{useUserLocation.ts}, abstraindo a complexidade dos sensores de geolocalização nativos.

\section{Visão de Dados}

O modelo de persistência baseia-se no paradigma Relacional para salvaguardar a consistência da FSM. Atualmente, o protótipo utiliza o \textit{SGBD SQLite} para validar o fluxo de ponta a ponta com \textit{optimistic locking}, mantendo compatibilidade arquitetural para migração para \textit{PostgreSQL} em escala de produção. O mapeamento é orquestrado via \textit{SQLAlchemy}, estruturando o esquema lógico em torno de três entidades fundamentais: \texttt{users}, para metadados de identificação; \texttt{pokemons}, para referências da PokeAPI e estado da entidade; e \texttt{adoptions}, uma tabela de junção que materializa o vínculo transacional e audita o ciclo de vida da adoção.

\section{Visão Física e de Implantação}

A infraestrutura distribui as cargas de processamento em três nós lógicos de implantação. A Camada Cliente consiste nos artefatos compilados via \textit{Expo}, que executam nativamente em dispositivos móveis consumindo o hardware do receptor GPS. A Camada de           
\chapter{Explicação de Trechos Importantes de Código}

A implementação do sistema fundamenta-se na integração entre a captura de sensores nativos no \textit{frontend} e a validação estrita de regras de negócio no \textit{backend}. Abaixo, detalham-se os componentes críticos que operacionalizam a telemetria, o gerenciamento de identidade e a validação espacial.

\section{Captura e Telemetria de Geolocalização}

A interface móvel abstrai a complexidade do hardware via \textit{hooks} customizados. O \texttt{useUserLocation.ts} implementa a observação passiva de coordenadas, integrando o \texttt{expo-location} para garantir a precisão necessária às validações da FSM. A escolha pela precisão \texttt{High} e um intervalo de \texttt{timeInterval: 5000} equilibra a responsividade do radar com o consumo de energia do dispositivo.

\begin{lstlisting}[style=estiloJS, caption={Hook de telemetria para captura de geolocalização}]
export const useUserLocation = () => {
  const [location, setLocation] = useState<Location.LocationObject | null>(null);

  useEffect(() => {
    (async () => {
      let { status } = await Location.requestForegroundPermissionsAsync();
      if (status !== 'granted') return;

      await Location.watchPositionAsync({
        accuracy: Location.Accuracy.High,
        timeInterval: 5000,
      }, (newLocation) => {
        setLocation(newLocation);
      });
    })();
  }, []);

  return location;
};
\end{lstlisting}

\section{Gerenciamento de Autenticação e Navegação}

O controle de sessão é consolidado no \texttt{AuthContext.tsx}, que atua como o \textit{Provider} central. A arquitetura evita a propagação excessiva de propriedades (\textit{prop-drilling}) ao injetar o estado do usuário diretamente na árvore de navegação. No \texttt{RootNavigator.tsx}, a condicional de renderização assegura que o estado de autenticação dite a composição da \textit{stack} de navegação, roteando o usuário dinamicamente entre as telas de \textit{login} e a aplicação funcional (\texttt{MainTabs.tsx}).

\begin{lstlisting}[style=estiloJS, caption={Navegação condicional baseada no contexto de autenticação}]
const RootNavigator = () => {
  const { user } = useContext(AuthContext);

  return (
    <NavigationContainer>
      {user ? <MainTabs /> : <LoginScreen />}
    </NavigationContainer>
  );
};
\end{lstlisting}

\section{Validação Espacial no Servidor}

Diferente do \textit{frontend}, que provê a telemetria, a inteligência do sistema reside no \texttt{spatial\_service.py}. A validação de proximidade para a adoção de um Pokémon não confia no dispositivo do usuário; o servidor recalcula a distância geodésica em relação à coordenada absoluta da entidade persistida no banco.

\begin{lstlisting}[style=estiloPython, caption={Validação de proximidade via Haversine no servidor}]
def is_within_distance(lat1, lon1, lat2, lon2, limit=50):
    dlat = radians(lat2 - lat1)
    dlon = radians(lon2 - lon1)
    a = sin(dlat/2)**2 + cos(radians(lat1)) * cos(radians(lat2)) * sin(dlon/2)**2
    c = 2 * atan2(sqrt(a), sqrt(1-a))
    distance = 6371000 * c  # Distancia em metros
    return distance <= limit
\end{lstlisting}

Este trecho é o núcleo da segurança da FSM. Ao executar esta validação isoladamente no servidor, o sistema blinda-se contra \textit{spoofing} de GPS, garantindo que o estado de \textit{ADOPTED} só seja atingido quando a proximidade física for matematicamente comprovada pela fórmula de Haversine.            
\chapter{Testes e Validação}

Este capítulo detalha as estratégias de garantia de qualidade (QA) e as metodologias de validação aplicadas para atestar a integridade, o desempenho e a segurança do ecossistema Poke-Adocao. O processo de homologação afasta-se de validações empíricas manuais em prol de uma esteira de testes automatizados com cobertura arquitetural contínua.

\section{Cenários de Teste e Automação}

A validação do fluxo transacional de ponta a ponta (\textit{End-to-End}) é assegurada por testes unitários e de integração rigorosos, divididos entre os dois pólos do monorrepositório, isolando as responsabilidades de estado e de interface.

\subsection{Validação do \textit{Backend}}
No ambiente de servidor, a infraestrutura de testes é orquestrada pela \textit{framework} \textit{Pytest}, acoplada à extensão \textit{pytest-asyncio} para suportar a resolução do \textit{event loop} do FastAPI. O banco de dados relacional é interceptado, e as transações de teste ocorrem em bases \textit{SQLite} em memória para evitar poluição de dados persistentes.
Os cenários homologados na base de código garantem as invariantes do domínio espacial e transacional:
\begin{itemize}
    \item \textbf{Motor Espacial (\texttt{test\_spatial\_service.py}):} Testes focados em atestar a precisão matemática da Fómula de Haversine e a delimitação dos vértices da \textit{Bounding Box} contra anomalias trigonométricas de latitude e longitude.
    \item \textbf{Máquina de Estados (\texttt{test\_adoption\_service.py}):} Validação da transição unidirecional da \textit{AdoptionStatus}, forçando falhas deliberadas como adoção concorrente da mesma entidade (\textit{race conditions}), violação da distância máxima de 50 metros e excedente do limite de inventário da \textit{party}.
    \item \textbf{Contratos de API (\texttt{test\_api.py} e \texttt{test\_auth.py}):} Verificação da emissão de \textit{tokens} criptográficos, proteção de rotas restritas e sanitização de requisições HTTP malformadas.
\end{itemize}

\subsection{Validação do \textit{Frontend}}
Na camada cliente, a integridade da aplicação móvel \textit{React Native} é aferida através do motor \textit{Jest} e das bibliotecas de dom \textit{React Native Testing Library}.
Os testes isolam o comportamento da interface do usuário da latência de rede:
\begin{itemize}
    \item \textbf{Comportamento de Hooks (\texttt{useUserLocation.test.ts}):} Simulação do ciclo de vida dos sensores do dispositivo móvel para assegurar que a telemetria geográfica é corretamente capturada, filtrada e exposta ao contexto da aplicação.
    \item \textbf{Renderização de Telas (\texttt{LoginScreen}, \texttt{PartyScreen}, \texttt{ProfileScreen}):} Atesta-se a montagem da árvore de componentes, garantindo que o inventário (\textit{Party}) não renderiza instâncias nulas e que o roteamento restringe utilizadores não autenticados.
\end{itemize}

\section{Visão do Usuário}

\begin{figure}[htb]
    \centering
    \caption{Tela de Login}
    \label{fig:login}
    \includegraphics[width=\linewidth, height=0.6\textheight, keepaspectratio]{figure/login.jpeg} 
    \fonte{Elaborado pelos autores (2026).}
\end{figure}

\begin{figure}[htb]
    \centering
    \caption{Tela do grupo de pokemons do usuário (party)}
    \label{fig:party}
    \includegraphics[width=\linewidth, height=0.6\textheight, keepaspectratio]{figure/party.jpeg} 
    \fonte{Elaborado pelos autores (2026).}
\end{figure}

\begin{figure}[htb]
    \centering
    \caption{Tela da lista de adoções vazia}
    \label{fig:lista_de_adocao}
    \includegraphics[width=\linewidth, height=0.6\textheight, keepaspectratio]{figure/lista_de_adocao.jpeg} 
    \fonte{Elaborado pelos autores (2026).}
\end{figure}

\begin{figure}[htb]
    \centering
    \caption{Tela da lista de adoções populada}
    \label{fig:lista_de_adocao_populada}
    \includegraphics[width=\linewidth, height=0.6\textheight, keepaspectratio]{figure/populada.jpeg} 
    \fonte{Elaborado pelos autores (2026).}
\end{figure}

\begin{figure}[htb]
    \centering
    \caption{Tela do Radar (Sem Pokemons ou jogadores perto)}
    \label{fig:radar}
    \includegraphics[width=\linewidth, height=0.6\textheight, keepaspectratio]{figure/radar (indisponivel).jpeg} 
    \fonte{Elaborado pelos autores (2026).}
\end{figure}

\begin{figure}[htb]
    \centering
    \caption{Tela do treinador (parte 1)}
    \label{fig:treinadorum}
    \includegraphics[width=\linewidth, height=0.6\textheight, keepaspectratio]{figure/treinador.jpeg} 
    \fonte{Elaborado pelos autores (2026).}
\end{figure}

\begin{figure}[htb]
    \centering
    \caption{Tela do treinador (parte 2)}
    \label{fig:treinadordois}
    \includegraphics[width=\linewidth, height=0.6\textheight, keepaspectratio]{figure/treinador (final).jpeg} 
    \fonte{Elaborado pelos autores (2026).}
\end{figure}           
\chapter{Conclusão}

O desenvolvimento do sistema Poke-Adocao permitiu a concretização da arquitetura de geolocalização e gerenciamento de estado proposta inicialmente. Os objetivos gerais e específicos delineados no Capítulo 3 foram integralmente alcançados, com a entrega de um sistema funcional que integra uma API robusta de \textit{backend} em Python e FastAPI a uma interface móvel multiplataforma em React Native. A validação das regras de negócio, assegurada pela Máquina de Estados Finitos (FSM) e pelo motor espacial baseado na fórmula de Haversine, demonstrou que a arquitetura proposta é capaz de garantir a integridade transacional mesmo sob condições de concorrência e restrições de proximidade física.

Do ponto de vista da engenharia de \textit{software}, os maiores gargalos técnicos enfrentados concentraram-se na mitigação da latência de processamento espacial e na garantia da consistência de dados em ambiente distribuído. A complexidade de calcular distâncias geodésicas em tempo real, aliada à necessidade de evitar condições de corrida (\textit{race conditions}) na adoção de entidades únicas, exigiu a implementação de uma camada de bloqueio otimista (\textit{optimistic locking}) no banco de dados. A transição para um modelo de monorrepositório foi determinante para superar a fricção de desalinhamento de contratos entre as camadas, permitindo que a equipe focasse no refinamento dos algoritmos de validação em vez de na gestão de dependências externas.

Quanto aos trabalhos futuros, o projeto mantém uma lista estratégica de dívidas técnicas e melhorias planejadas para a próxima fase de maturação. O plano de expansão inclui:
\begin{itemize}
    \item \textbf{Escalabilidade Geográfica:} Implementação de índices espaciais nativos (\textit{PostGIS}) para substituir a pré-filtragem por \textit{Bounding Box} manual, elevando a performance das consultas para cenários de alta densidade de usuários.
    \item \textbf{Dívida Técnica de \textit{Frontend}:} Migração dos componentes de interface para um sistema de temas totalmente tokenizado, aumentando a manutenibilidade do design.
    \item \textbf{Novas Funcionalidades:} Desenvolvimento de um motor de \textit{caching} distribuído (Redis) para persistir temporariamente as telemetrias de localização dos usuários e reduzir o custo de escrita direta no banco relacional.
\end{itemize}

Em suma, o projeto Poke-Adocao demonstra a viabilidade de construir aplicações de \textit{Location-Based Services} complexas utilizando um \textit{stack} tecnológico moderno e rigoroso, estabelecendo uma fundação sólida para futuras iterações e melhorias de escalabilidade de engenharia.   
\include{capitulos/05-trechos-codigo}       
@book{fowler2002patterns,
  title={Patterns of Enterprise Application Architecture},
  author={Fowler, Martin},
  year={2002},
  publisher={Addison-Wesley Longman},
  address={Boston}
}

@article{sinnott1984virtues,
  title={Virtues of the Haversine},
  author={Sinnott, Roger W},
  journal={Sky and Telescope},
  volume={68},
  number={2},
  pages={159},
  year={1984}
}

@phdthesis{fielding2000architectural,
  title={Architectural styles and the design of network-based software architectures},
  author={Fielding, Roy Thomas},
  year={2000},
  school={University of California, Irvine}
}

@book{gurnani2020developing,
  title={Developing Applications with React Native},
  author={Gurnani, Akshat},
  year={2020},
  publisher={Packt Publishing Ltd},
  address={Birmingham}
}

@book{sommerville2011software,
  title={Software Engineering},
  author={Sommerville, Ian},
  year={2011},
  publisher={Addison-Wesley},
  address={Boston}
}

% ==============================================================================
% MANUAL DE OPERAÇÃO: GERENCIAMENTO DE REFERÊNCIAS (BIBTEX + ABNTEX2)
% ==============================================================================
% 
% 1. O CONCEITO DO ARQUIVO .BIB
% No LaTeX, você não digita a lista de referências manualmente no final do texto.
% O arquivo 'referencias.bib' atua como um banco de dados relacional. Você
% cadastra os metadados das fontes nele e faz "consultas" no texto principal.
% Regra de Ouro: Uma fonte cadastrada no .bib SÓ aparecerá no PDF final se
% for explicitamente citada em algum lugar do texto.
%
% 2. COMO CADASTRAR UMA FONTE (No arquivo referencias.bib)
% Cada entrada possui um "tipo" (@book, @article, @misc) e uma "chave" única
% que você inventa para chamá-la depois.
%
% Exemplo de Livro Técnico:
% @book{sommerville2011,
%   author    = {Sommerville, Ian},
%   title     = {Engenharia de software},
%   publisher = {Pearson Prentice Hall},
%   year      = {2011},
%   address   = {São Paulo}
% }
%
% Exemplo de Site, Documentação ou Repositório:
% @misc{fastapi2026,
%   author        = {Ramírez, Sebastián},
%   title         = {FastAPI Documentation},
%   year          = {2026},
%   url           = {https://fastapi.tiangolo.com/},
%   urlaccessdate = {8 jun. 2026}
% }
% Nota: Autores corporativos (ex: Meta, Google, W3C) devem ser envolvidos 
% em chaves duplas no arquivo .bib para que o LaTeX não inverta os nomes. 
% Ex: author = {{Meta Platforms}}
%
% 3. COMO CHAMAR A FONTE NO TEXTO (Nos arquivos .tex)
% O pacote abntex2cite fornece comandos rigorosos para adequação à NBR 10520.
%
% -> Citação Indireta (Ao final da frase, em maiúsculas):
% Comando: \cite{chave}
% Exemplo: A arquitetura baseada em eventos reduz o acoplamento \cite{sommerville2011}.
% PDF renderizado: A arquitetura baseada em eventos reduz o acoplamento (SOMMERVILLE, 2011).
%
% -> Citação Direta/Narrativa (Inserida no fluxo do texto):
% Comando: \citeonline{chave}
% Exemplo: Segundo \citeonline{sommerville2011}, a engenharia de requisitos é vital.
% PDF renderizado: Segundo Sommerville (2011), a engenharia de requisitos é vital.
%
% -> Inserção Forçada (Sem citação textual):
% Comando: \nocite{chave}
% Uso: Insere a fonte na bibliografia final sem precisar colocar um (AUTOR, Ano) no texto.
% Evite usar isso em trabalhos acadêmicos para burlar a fundamentação.
% ==============================================================================     
% ------------------------------------------------------------------------
% Configuração e Arquitetura do Template LaTeX
% Desenvolvido por: Aurora Drumond Costa Magalhães
% Ano: 2026
% ------------------------------------------------------------------------

\documentclass[
    12pt,
    openright,
    oneside,
    a4paper,
    english,
    brazil
    ]{abntex2}

% --- PACOTES GERAIS ---
\usepackage{lmodern}
\usepackage[T1]{fontenc}
\usepackage[utf8]{inputenc}
\usepackage{indentfirst}
\usepackage{graphicx}
\usepackage{float}
\usepackage[alf]{abntex2cite}

% --- CONFIGURAÇÃO DO AMBIENTE QUADRO ---
\newcommand{\quadroname}{Quadro}
\newcommand{\listofquadrosname}{Lista de quadros}

% Utilizando a sintaxe estrita do pacote 'float' (\newfloat{nome}{posicionamento}{extensão}[contador])
\floatstyle{plaintop} % Garante que a legenda fique acima do quadro
\newfloat{quadro}{htb}{loq}[chapter]
\floatname{quadro}{\quadroname}

% Cria a lista de quadros associada ao sumário
\newlistof{listofquadros}{loq}{\listofquadrosname}
\newlistentry{quadro}{loq}{0}

% Define a numeração contínua (evita Quadro 1.1, Quadro 1.2)
\counterwithout{quadro}{chapter}

% Configura a formatação de exibição da Lista de Quadros no PDF
\renewcommand{\cftquadroname}{\quadroname\space}
\renewcommand*{\cftquadroaftersnum}{\hfill--\hfill}

% --- IMPORTA SEUS ESTILOS PERSONALIZADOS ---
% Substitua pelos caminhos reais dos seus arquivos de estilo (.sty), caso possua.
\usepackage{styles/code-style} 
\usepackage{styles/mylib-style}

% --- DADOS DA CAPA E FOLHA DE ROSTO ---
\titulo{Poke-Adocao: Sistema de Adoção de Pokémon via Geolocalização}
\autor{Aurora Drumond Costa Magalhães \and Ana Clara de Souza \and Kayke Wellington}
\local{Timóteo - MG}
\data{2026}

\instituicao{%
  Centro Universitário Católica do Leste de Minas Gerais
  \par
  Coronel Fabriciano}

\tipotrabalho{Documentação de Arquitetura de Software}

\preambulo{Relatório técnico apresentado como documentação oficial da arquitetura do sistema Poke-Adocao, detalhando as especificações de engenharia do serviço de backend e da aplicação móvel.}

% --- INÍCIO DO DOCUMENTO ---
\begin{document}

% Elementos Pré-Textuais (Obrigatórios segundo a ABNT)
\imprimircapa
\imprimirfolhaderosto

% --- RESUMOS ---
\setlength{\absparsep}{18pt}

\begin{resumo}
O projeto Poke-Adocao consiste num sistema de adoção de Pokémon que utiliza geolocalização para interligar os utilizadores a criaturas virtuais no espaço físico. A arquitetura baseia-se num monorepo que integra uma API RESTful em Python com FastAPI no backend, e uma aplicação móvel multiplataforma construída em React Native e Expo no frontend. O núcleo do processo de adoção é gerido por uma Máquina de Estados Finitos (FSM). Atualmente operando de forma binária para garantir a integridade do ciclo de vida restrito a um limite rigoroso de 50 metros de distância, a FSM demonstra a fundação arquitetural necessária para suportar, futuramente, estados intermediários mais dinâmicos e contratos \textit{player-to-player} (P2P) complexos, como transações pendentes de aprovação. Para viabilizar este mecanismo, o sistema implementa um motor espacial que calcula distâncias geográficas utilizando a fórmula de Haversine diretamente na camada da aplicação, processo previamente otimizado por uma filtragem através de \textit{Bounding Box}. A persistência de dados é assegurada de forma leve e contida por um banco de dados relacional SQLite, escolhido estrategicamente para validar as transações de \textit{optimistic locking} e o fluxo de ponta a ponta do protótipo sem a sobrecarga de configuração de infraestrutura externa. O presente documento detalha as decisões de engenharia, a infraestrutura e os requisitos técnicos que sustentam a aplicação.

\vspace{\onelineskip}
\noindent
\textbf{Palavras-chave:} Geolocalização. React Native. FastAPI. Máquina de Estados Finitos. Haversine. SQLite.
\end{resumo}

\begin{resumo}[Abstract]
\begin{otherlanguage*}{english}
The Poke-Adocao project is a Pokemon adoption system that leverages geolocation to connect users with virtual creatures in the physical world. The architecture is built upon a monorepo integrating a RESTful API developed in Python with FastAPI for the backend, and a cross-platform mobile application built with React Native and Expo for the frontend. The core adoption process is managed by a Finite State Machine (FSM). Currently operating in a binary manner to enforce the lifecycle under a strict 50-meter geographic threshold, the FSM demonstrates the architectural foundation necessary to support complex intermediate states and dynamic player-to-player (P2P) contracts in the future, such as pending approvals. To enable this mechanism, the system implements a spatial engine that computes geographic distances using the Haversine formula at the application layer, optimized by a Bounding Box pre-filtering process. Data persistence is handled in a lightweight and contained manner by a relational SQLite database, strategically chosen to validate optimistic locking transactions and the end-to-end flow of the prototype without the overhead of external infrastructure configuration. This document details the engineering decisions, infrastructure, and technical requirements underlying the application.

\vspace{\onelineskip}
\noindent
\textbf{Keywords:} Geolocation. React Native. FastAPI. Finite State Machine. Haversine. SQLite.
\end{otherlanguage*}
\end{resumo}

% --- LISTAS ---
% Descomente as listas (removendo o sinal '%') que forem necessárias para o seu trabalho de acordo com os elementos presentes.

% Lista de Figuras
\pdfbookmark[0]{\listfigurename}{lof}
\listoffigures*
\cleardoublepage

% Lista de Códigos / Trechos de Software
% \pdfbookmark[0]{\lstlistlistingname}{lol}
% \lstlistoflistings
% \cleardoublepage

% Lista de Quadros
\pdfbookmark[0]{Lista de quadros}{loq}
\listofquadros*
\cleardoublepage

% Sumário (Obrigatório)
\pdfbookmark[0]{\contentsname}{toc}
\tableofcontents
\cleardoublepage

% --------------------------------------------------------
% ELEMENTOS TEXTUAIS
% --------------------------------------------------------
\textual

% --- ESTRUTURAÇÃO EM MÚLTIPLOS ARQUIVOS ---
% A chamada \include insere o conteúdo de arquivos externos. Isso permite modularizar o 
% documento separando os tópicos exigidos no relatório técnico em arquivos .tex 
% distintos dentro do diretório 'capitulos/'.
% Não é necessário adicionar a extensão .tex ao usar o comando \include.

\chapter{Design e Arquitetura}

O padrão arquitetural macro adotado pelo sistema Poke-Adocao fundamenta-se no modelo Cliente-Servidor Desacoplado, operacionalizado sob a estrutura de um monorrepositório. Essa abordagem centraliza a governança do código-fonte e dos contratos de comunicação, permitindo o isolamento completo entre a aplicação cliente móvel e a API de serviços, que interagem estritamente por meio de protocolos de rede \textit{stateless}.

\section{Visão Lógica e de Componentes}

A organização interna do ecossistema é bifurcada entre as camadas funcionais de \textit{backend} e \textit{frontend}, estruturadas de forma modular para assegurar a separação de conceitos \cite{sommerville2011software}. 

No \textit{backend}, a API segue uma Arquitetura em Camadas dirigida por serviços. A camada de roteamento e entrada, situada no arquivo \texttt{main.py}, atua como a interface de exposição dos \textit{endpoints} HTTP, gerenciando \textit{middlewares} e a orquestração de controladores assíncronos. A validação dos dados é centralizada no \texttt{schemas.py}, que utiliza o \textit{Pydantic} para garantir a tipagem estática e serialização de todas as requisições. As regras de negócio operam em uma camada dedicada composta por serviços como \texttt{adoption\_service.py} e \texttt{spatial\_service.py}, que isolam a lógica da Máquina de Estados Finitos (FSM) e as validações trigonométricas, enquanto a persistência é orquestrada através dos modelos declarativos em \texttt{models.py} e \texttt{database.py}, gerindo o ciclo de vida das sessões transacionais.

No \textit{frontend}, a interface adota uma arquitetura baseada em componentes funcionais alimentados por injeção de estado via \textit{hooks}. O \texttt{RootNavigator.tsx} gerencia o fluxo global de autenticação condicional, chaveando o utilizador entre a tela de \textit{login} e a navegação principal estruturada no \texttt{MainTabs.tsx}. A gestão de estado é centralizada nativamente no \texttt{AuthContext.tsx}, encapsulando a persistência de \textit{tokens} e simplificando a telemetria, que é isolada no \textit{hook} reativo \texttt{useUserLocation.ts}, abstraindo a complexidade dos sensores de geolocalização nativos.

\section{Visão de Dados}

O modelo de persistência baseia-se no paradigma Relacional para salvaguardar a consistência da FSM. Atualmente, o protótipo utiliza o \textit{SGBD SQLite} para validar o fluxo de ponta a ponta com \textit{optimistic locking}, mantendo compatibilidade arquitetural para migração para \textit{PostgreSQL} em escala de produção. O mapeamento é orquestrado via \textit{SQLAlchemy}, estruturando o esquema lógico em torno de três entidades fundamentais: \texttt{users}, para metadados de identificação; \texttt{pokemons}, para referências da PokeAPI e estado da entidade; e \texttt{adoptions}, uma tabela de junção que materializa o vínculo transacional e audita o ciclo de vida da adoção.

\section{Visão Física e de Implantação}

A infraestrutura distribui as cargas de processamento em três nós lógicos de implantação. A Camada Cliente consiste nos artefatos compilados via \textit{Expo}, que executam nativamente em dispositivos móveis consumindo o hardware do receptor GPS. A Camada de           
\chapter{Explicação de Trechos Importantes de Código}

A implementação do sistema fundamenta-se na integração entre a captura de sensores nativos no \textit{frontend} e a validação estrita de regras de negócio no \textit{backend}. Abaixo, detalham-se os componentes críticos que operacionalizam a telemetria, o gerenciamento de identidade e a validação espacial.

\section{Captura e Telemetria de Geolocalização}

A interface móvel abstrai a complexidade do hardware via \textit{hooks} customizados. O \texttt{useUserLocation.ts} implementa a observação passiva de coordenadas, integrando o \texttt{expo-location} para garantir a precisão necessária às validações da FSM. A escolha pela precisão \texttt{High} e um intervalo de \texttt{timeInterval: 5000} equilibra a responsividade do radar com o consumo de energia do dispositivo.

\begin{lstlisting}[style=estiloJS, caption={Hook de telemetria para captura de geolocalização}]
export const useUserLocation = () => {
  const [location, setLocation] = useState<Location.LocationObject | null>(null);

  useEffect(() => {
    (async () => {
      let { status } = await Location.requestForegroundPermissionsAsync();
      if (status !== 'granted') return;

      await Location.watchPositionAsync({
        accuracy: Location.Accuracy.High,
        timeInterval: 5000,
      }, (newLocation) => {
        setLocation(newLocation);
      });
    })();
  }, []);

  return location;
};
\end{lstlisting}

\section{Gerenciamento de Autenticação e Navegação}

O controle de sessão é consolidado no \texttt{AuthContext.tsx}, que atua como o \textit{Provider} central. A arquitetura evita a propagação excessiva de propriedades (\textit{prop-drilling}) ao injetar o estado do usuário diretamente na árvore de navegação. No \texttt{RootNavigator.tsx}, a condicional de renderização assegura que o estado de autenticação dite a composição da \textit{stack} de navegação, roteando o usuário dinamicamente entre as telas de \textit{login} e a aplicação funcional (\texttt{MainTabs.tsx}).

\begin{lstlisting}[style=estiloJS, caption={Navegação condicional baseada no contexto de autenticação}]
const RootNavigator = () => {
  const { user } = useContext(AuthContext);

  return (
    <NavigationContainer>
      {user ? <MainTabs /> : <LoginScreen />}
    </NavigationContainer>
  );
};
\end{lstlisting}

\section{Validação Espacial no Servidor}

Diferente do \textit{frontend}, que provê a telemetria, a inteligência do sistema reside no \texttt{spatial\_service.py}. A validação de proximidade para a adoção de um Pokémon não confia no dispositivo do usuário; o servidor recalcula a distância geodésica em relação à coordenada absoluta da entidade persistida no banco.

\begin{lstlisting}[style=estiloPython, caption={Validação de proximidade via Haversine no servidor}]
def is_within_distance(lat1, lon1, lat2, lon2, limit=50):
    dlat = radians(lat2 - lat1)
    dlon = radians(lon2 - lon1)
    a = sin(dlat/2)**2 + cos(radians(lat1)) * cos(radians(lat2)) * sin(dlon/2)**2
    c = 2 * atan2(sqrt(a), sqrt(1-a))
    distance = 6371000 * c  # Distancia em metros
    return distance <= limit
\end{lstlisting}

Este trecho é o núcleo da segurança da FSM. Ao executar esta validação isoladamente no servidor, o sistema blinda-se contra \textit{spoofing} de GPS, garantindo que o estado de \textit{ADOPTED} só seja atingido quando a proximidade física for matematicamente comprovada pela fórmula de Haversine.            
\chapter{Testes e Validação}

Este capítulo detalha as estratégias de garantia de qualidade (QA) e as metodologias de validação aplicadas para atestar a integridade, o desempenho e a segurança do ecossistema Poke-Adocao. O processo de homologação afasta-se de validações empíricas manuais em prol de uma esteira de testes automatizados com cobertura arquitetural contínua.

\section{Cenários de Teste e Automação}

A validação do fluxo transacional de ponta a ponta (\textit{End-to-End}) é assegurada por testes unitários e de integração rigorosos, divididos entre os dois pólos do monorrepositório, isolando as responsabilidades de estado e de interface.

\subsection{Validação do \textit{Backend}}
No ambiente de servidor, a infraestrutura de testes é orquestrada pela \textit{framework} \textit{Pytest}, acoplada à extensão \textit{pytest-asyncio} para suportar a resolução do \textit{event loop} do FastAPI. O banco de dados relacional é interceptado, e as transações de teste ocorrem em bases \textit{SQLite} em memória para evitar poluição de dados persistentes.
Os cenários homologados na base de código garantem as invariantes do domínio espacial e transacional:
\begin{itemize}
    \item \textbf{Motor Espacial (\texttt{test\_spatial\_service.py}):} Testes focados em atestar a precisão matemática da Fómula de Haversine e a delimitação dos vértices da \textit{Bounding Box} contra anomalias trigonométricas de latitude e longitude.
    \item \textbf{Máquina de Estados (\texttt{test\_adoption\_service.py}):} Validação da transição unidirecional da \textit{AdoptionStatus}, forçando falhas deliberadas como adoção concorrente da mesma entidade (\textit{race conditions}), violação da distância máxima de 50 metros e excedente do limite de inventário da \textit{party}.
    \item \textbf{Contratos de API (\texttt{test\_api.py} e \texttt{test\_auth.py}):} Verificação da emissão de \textit{tokens} criptográficos, proteção de rotas restritas e sanitização de requisições HTTP malformadas.
\end{itemize}

\subsection{Validação do \textit{Frontend}}
Na camada cliente, a integridade da aplicação móvel \textit{React Native} é aferida através do motor \textit{Jest} e das bibliotecas de dom \textit{React Native Testing Library}.
Os testes isolam o comportamento da interface do usuário da latência de rede:
\begin{itemize}
    \item \textbf{Comportamento de Hooks (\texttt{useUserLocation.test.ts}):} Simulação do ciclo de vida dos sensores do dispositivo móvel para assegurar que a telemetria geográfica é corretamente capturada, filtrada e exposta ao contexto da aplicação.
    \item \textbf{Renderização de Telas (\texttt{LoginScreen}, \texttt{PartyScreen}, \texttt{ProfileScreen}):} Atesta-se a montagem da árvore de componentes, garantindo que o inventário (\textit{Party}) não renderiza instâncias nulas e que o roteamento restringe utilizadores não autenticados.
\end{itemize}

\section{Visão do Usuário}

\begin{figure}[htb]
    \centering
    \caption{Tela de Login}
    \label{fig:login}
    \includegraphics[width=\linewidth, height=0.6\textheight, keepaspectratio]{figure/login.jpeg} 
    \fonte{Elaborado pelos autores (2026).}
\end{figure}

\begin{figure}[htb]
    \centering
    \caption{Tela do grupo de pokemons do usuário (party)}
    \label{fig:party}
    \includegraphics[width=\linewidth, height=0.6\textheight, keepaspectratio]{figure/party.jpeg} 
    \fonte{Elaborado pelos autores (2026).}
\end{figure}

\begin{figure}[htb]
    \centering
    \caption{Tela da lista de adoções vazia}
    \label{fig:lista_de_adocao}
    \includegraphics[width=\linewidth, height=0.6\textheight, keepaspectratio]{figure/lista_de_adocao.jpeg} 
    \fonte{Elaborado pelos autores (2026).}
\end{figure}

\begin{figure}[htb]
    \centering
    \caption{Tela da lista de adoções populada}
    \label{fig:lista_de_adocao_populada}
    \includegraphics[width=\linewidth, height=0.6\textheight, keepaspectratio]{figure/populada.jpeg} 
    \fonte{Elaborado pelos autores (2026).}
\end{figure}

\begin{figure}[htb]
    \centering
    \caption{Tela do Radar (Sem Pokemons ou jogadores perto)}
    \label{fig:radar}
    \includegraphics[width=\linewidth, height=0.6\textheight, keepaspectratio]{figure/radar (indisponivel).jpeg} 
    \fonte{Elaborado pelos autores (2026).}
\end{figure}

\begin{figure}[htb]
    \centering
    \caption{Tela do treinador (parte 1)}
    \label{fig:treinadorum}
    \includegraphics[width=\linewidth, height=0.6\textheight, keepaspectratio]{figure/treinador.jpeg} 
    \fonte{Elaborado pelos autores (2026).}
\end{figure}

\begin{figure}[htb]
    \centering
    \caption{Tela do treinador (parte 2)}
    \label{fig:treinadordois}
    \includegraphics[width=\linewidth, height=0.6\textheight, keepaspectratio]{figure/treinador (final).jpeg} 
    \fonte{Elaborado pelos autores (2026).}
\end{figure}           
\chapter{Conclusão}

O desenvolvimento do sistema Poke-Adocao permitiu a concretização da arquitetura de geolocalização e gerenciamento de estado proposta inicialmente. Os objetivos gerais e específicos delineados no Capítulo 3 foram integralmente alcançados, com a entrega de um sistema funcional que integra uma API robusta de \textit{backend} em Python e FastAPI a uma interface móvel multiplataforma em React Native. A validação das regras de negócio, assegurada pela Máquina de Estados Finitos (FSM) e pelo motor espacial baseado na fórmula de Haversine, demonstrou que a arquitetura proposta é capaz de garantir a integridade transacional mesmo sob condições de concorrência e restrições de proximidade física.

Do ponto de vista da engenharia de \textit{software}, os maiores gargalos técnicos enfrentados concentraram-se na mitigação da latência de processamento espacial e na garantia da consistência de dados em ambiente distribuído. A complexidade de calcular distâncias geodésicas em tempo real, aliada à necessidade de evitar condições de corrida (\textit{race conditions}) na adoção de entidades únicas, exigiu a implementação de uma camada de bloqueio otimista (\textit{optimistic locking}) no banco de dados. A transição para um modelo de monorrepositório foi determinante para superar a fricção de desalinhamento de contratos entre as camadas, permitindo que a equipe focasse no refinamento dos algoritmos de validação em vez de na gestão de dependências externas.

Quanto aos trabalhos futuros, o projeto mantém uma lista estratégica de dívidas técnicas e melhorias planejadas para a próxima fase de maturação. O plano de expansão inclui:
\begin{itemize}
    \item \textbf{Escalabilidade Geográfica:} Implementação de índices espaciais nativos (\textit{PostGIS}) para substituir a pré-filtragem por \textit{Bounding Box} manual, elevando a performance das consultas para cenários de alta densidade de usuários.
    \item \textbf{Dívida Técnica de \textit{Frontend}:} Migração dos componentes de interface para um sistema de temas totalmente tokenizado, aumentando a manutenibilidade do design.
    \item \textbf{Novas Funcionalidades:} Desenvolvimento de um motor de \textit{caching} distribuído (Redis) para persistir temporariamente as telemetrias de localização dos usuários e reduzir o custo de escrita direta no banco relacional.
\end{itemize}

Em suma, o projeto Poke-Adocao demonstra a viabilidade de construir aplicações de \textit{Location-Based Services} complexas utilizando um \textit{stack} tecnológico moderno e rigoroso, estabelecendo uma fundação sólida para futuras iterações e melhorias de escalabilidade de engenharia.   
\include{capitulos/05-trechos-codigo}       
@book{fowler2002patterns,
  title={Patterns of Enterprise Application Architecture},
  author={Fowler, Martin},
  year={2002},
  publisher={Addison-Wesley Longman},
  address={Boston}
}

@article{sinnott1984virtues,
  title={Virtues of the Haversine},
  author={Sinnott, Roger W},
  journal={Sky and Telescope},
  volume={68},
  number={2},
  pages={159},
  year={1984}
}

@phdthesis{fielding2000architectural,
  title={Architectural styles and the design of network-based software architectures},
  author={Fielding, Roy Thomas},
  year={2000},
  school={University of California, Irvine}
}

@book{gurnani2020developing,
  title={Developing Applications with React Native},
  author={Gurnani, Akshat},
  year={2020},
  publisher={Packt Publishing Ltd},
  address={Birmingham}
}

@book{sommerville2011software,
  title={Software Engineering},
  author={Sommerville, Ian},
  year={2011},
  publisher={Addison-Wesley},
  address={Boston}
}

% ==============================================================================
% MANUAL DE OPERAÇÃO: GERENCIAMENTO DE REFERÊNCIAS (BIBTEX + ABNTEX2)
% ==============================================================================
% 
% 1. O CONCEITO DO ARQUIVO .BIB
% No LaTeX, você não digita a lista de referências manualmente no final do texto.
% O arquivo 'referencias.bib' atua como um banco de dados relacional. Você
% cadastra os metadados das fontes nele e faz "consultas" no texto principal.
% Regra de Ouro: Uma fonte cadastrada no .bib SÓ aparecerá no PDF final se
% for explicitamente citada em algum lugar do texto.
%
% 2. COMO CADASTRAR UMA FONTE (No arquivo referencias.bib)
% Cada entrada possui um "tipo" (@book, @article, @misc) e uma "chave" única
% que você inventa para chamá-la depois.
%
% Exemplo de Livro Técnico:
% @book{sommerville2011,
%   author    = {Sommerville, Ian},
%   title     = {Engenharia de software},
%   publisher = {Pearson Prentice Hall},
%   year      = {2011},
%   address   = {São Paulo}
% }
%
% Exemplo de Site, Documentação ou Repositório:
% @misc{fastapi2026,
%   author        = {Ramírez, Sebastián},
%   title         = {FastAPI Documentation},
%   year          = {2026},
%   url           = {https://fastapi.tiangolo.com/},
%   urlaccessdate = {8 jun. 2026}
% }
% Nota: Autores corporativos (ex: Meta, Google, W3C) devem ser envolvidos 
% em chaves duplas no arquivo .bib para que o LaTeX não inverta os nomes. 
% Ex: author = {{Meta Platforms}}
%
% 3. COMO CHAMAR A FONTE NO TEXTO (Nos arquivos .tex)
% O pacote abntex2cite fornece comandos rigorosos para adequação à NBR 10520.
%
% -> Citação Indireta (Ao final da frase, em maiúsculas):
% Comando: \cite{chave}
% Exemplo: A arquitetura baseada em eventos reduz o acoplamento \cite{sommerville2011}.
% PDF renderizado: A arquitetura baseada em eventos reduz o acoplamento (SOMMERVILLE, 2011).
%
% -> Citação Direta/Narrativa (Inserida no fluxo do texto):
% Comando: \citeonline{chave}
% Exemplo: Segundo \citeonline{sommerville2011}, a engenharia de requisitos é vital.
% PDF renderizado: Segundo Sommerville (2011), a engenharia de requisitos é vital.
%
% -> Inserção Forçada (Sem citação textual):
% Comando: \nocite{chave}
% Uso: Insere a fonte na bibliografia final sem precisar colocar um (AUTOR, Ano) no texto.
% Evite usar isso em trabalhos acadêmicos para burlar a fundamentação.
% ==============================================================================     
% ------------------------------------------------------------------------
% Configuração e Arquitetura do Template LaTeX
% Desenvolvido por: Aurora Drumond Costa Magalhães
% Ano: 2026
% ------------------------------------------------------------------------

\documentclass[
    12pt,
    openright,
    oneside,
    a4paper,
    english,
    brazil
    ]{abntex2}

% --- PACOTES GERAIS ---
\usepackage{lmodern}
\usepackage[T1]{fontenc}
\usepackage[utf8]{inputenc}
\usepackage{indentfirst}
\usepackage{graphicx}
\usepackage{float}
\usepackage[alf]{abntex2cite}

% --- CONFIGURAÇÃO DO AMBIENTE QUADRO ---
\newcommand{\quadroname}{Quadro}
\newcommand{\listofquadrosname}{Lista de quadros}

% Utilizando a sintaxe estrita do pacote 'float' (\newfloat{nome}{posicionamento}{extensão}[contador])
\floatstyle{plaintop} % Garante que a legenda fique acima do quadro
\newfloat{quadro}{htb}{loq}[chapter]
\floatname{quadro}{\quadroname}

% Cria a lista de quadros associada ao sumário
\newlistof{listofquadros}{loq}{\listofquadrosname}
\newlistentry{quadro}{loq}{0}

% Define a numeração contínua (evita Quadro 1.1, Quadro 1.2)
\counterwithout{quadro}{chapter}

% Configura a formatação de exibição da Lista de Quadros no PDF
\renewcommand{\cftquadroname}{\quadroname\space}
\renewcommand*{\cftquadroaftersnum}{\hfill--\hfill}

% --- IMPORTA SEUS ESTILOS PERSONALIZADOS ---
% Substitua pelos caminhos reais dos seus arquivos de estilo (.sty), caso possua.
\usepackage{styles/code-style} 
\usepackage{styles/mylib-style}

% --- DADOS DA CAPA E FOLHA DE ROSTO ---
\titulo{Poke-Adocao: Sistema de Adoção de Pokémon via Geolocalização}
\autor{Aurora Drumond Costa Magalhães \and Ana Clara de Souza \and Kayke Wellington}
\local{Timóteo - MG}
\data{2026}

\instituicao{%
  Centro Universitário Católica do Leste de Minas Gerais
  \par
  Coronel Fabriciano}

\tipotrabalho{Documentação de Arquitetura de Software}

\preambulo{Relatório técnico apresentado como documentação oficial da arquitetura do sistema Poke-Adocao, detalhando as especificações de engenharia do serviço de backend e da aplicação móvel.}

% --- INÍCIO DO DOCUMENTO ---
\begin{document}

% Elementos Pré-Textuais (Obrigatórios segundo a ABNT)
\imprimircapa
\imprimirfolhaderosto

% --- RESUMOS ---
\setlength{\absparsep}{18pt}

\begin{resumo}
O projeto Poke-Adocao consiste num sistema de adoção de Pokémon que utiliza geolocalização para interligar os utilizadores a criaturas virtuais no espaço físico. A arquitetura baseia-se num monorepo que integra uma API RESTful em Python com FastAPI no backend, e uma aplicação móvel multiplataforma construída em React Native e Expo no frontend. O núcleo do processo de adoção é gerido por uma Máquina de Estados Finitos (FSM). Atualmente operando de forma binária para garantir a integridade do ciclo de vida restrito a um limite rigoroso de 50 metros de distância, a FSM demonstra a fundação arquitetural necessária para suportar, futuramente, estados intermediários mais dinâmicos e contratos \textit{player-to-player} (P2P) complexos, como transações pendentes de aprovação. Para viabilizar este mecanismo, o sistema implementa um motor espacial que calcula distâncias geográficas utilizando a fórmula de Haversine diretamente na camada da aplicação, processo previamente otimizado por uma filtragem através de \textit{Bounding Box}. A persistência de dados é assegurada de forma leve e contida por um banco de dados relacional SQLite, escolhido estrategicamente para validar as transações de \textit{optimistic locking} e o fluxo de ponta a ponta do protótipo sem a sobrecarga de configuração de infraestrutura externa. O presente documento detalha as decisões de engenharia, a infraestrutura e os requisitos técnicos que sustentam a aplicação.

\vspace{\onelineskip}
\noindent
\textbf{Palavras-chave:} Geolocalização. React Native. FastAPI. Máquina de Estados Finitos. Haversine. SQLite.
\end{resumo}

\begin{resumo}[Abstract]
\begin{otherlanguage*}{english}
The Poke-Adocao project is a Pokemon adoption system that leverages geolocation to connect users with virtual creatures in the physical world. The architecture is built upon a monorepo integrating a RESTful API developed in Python with FastAPI for the backend, and a cross-platform mobile application built with React Native and Expo for the frontend. The core adoption process is managed by a Finite State Machine (FSM). Currently operating in a binary manner to enforce the lifecycle under a strict 50-meter geographic threshold, the FSM demonstrates the architectural foundation necessary to support complex intermediate states and dynamic player-to-player (P2P) contracts in the future, such as pending approvals. To enable this mechanism, the system implements a spatial engine that computes geographic distances using the Haversine formula at the application layer, optimized by a Bounding Box pre-filtering process. Data persistence is handled in a lightweight and contained manner by a relational SQLite database, strategically chosen to validate optimistic locking transactions and the end-to-end flow of the prototype without the overhead of external infrastructure configuration. This document details the engineering decisions, infrastructure, and technical requirements underlying the application.

\vspace{\onelineskip}
\noindent
\textbf{Keywords:} Geolocation. React Native. FastAPI. Finite State Machine. Haversine. SQLite.
\end{otherlanguage*}
\end{resumo}

% --- LISTAS ---
% Descomente as listas (removendo o sinal '%') que forem necessárias para o seu trabalho de acordo com os elementos presentes.

% Lista de Figuras
\pdfbookmark[0]{\listfigurename}{lof}
\listoffigures*
\cleardoublepage

% Lista de Códigos / Trechos de Software
% \pdfbookmark[0]{\lstlistlistingname}{lol}
% \lstlistoflistings
% \cleardoublepage

% Lista de Quadros
\pdfbookmark[0]{Lista de quadros}{loq}
\listofquadros*
\cleardoublepage

% Sumário (Obrigatório)
\pdfbookmark[0]{\contentsname}{toc}
\tableofcontents
\cleardoublepage

% --------------------------------------------------------
% ELEMENTOS TEXTUAIS
% --------------------------------------------------------
\textual

% --- ESTRUTURAÇÃO EM MÚLTIPLOS ARQUIVOS ---
% A chamada \include insere o conteúdo de arquivos externos. Isso permite modularizar o 
% documento separando os tópicos exigidos no relatório técnico em arquivos .tex 
% distintos dentro do diretório 'capitulos/'.
% Não é necessário adicionar a extensão .tex ao usar o comando \include.

\include{capitulos/01-introducao}           
\include{capitulos/02-objetivos}            
\include{capitulos/03-requisitos}           
\include{capitulos/04-design-arquitetura}   
\include{capitulos/05-trechos-codigo}       
\include{capitulos/06-testes-validacao}     
\include{capitulos/07-conclusao}            

% --------------------------------------------------------
% ELEMENTOS PÓS-TEXTUAIS
% --------------------------------------------------------
\postextual

% --- REFERÊNCIAS BIBLIOGRÁFICAS ---
% Elemento obrigatório. Puxa as citações do arquivo referencias.bib formatando nas normas da ABNT (NBR 6023).
\bibliographystyle{abntex2-alf} % Define o estilo ABNT (Autor-Data)
\bibliography{references}

% --- APÊNDICES ---
% Elemento opcional. Documentos criados pelo próprio autor para complementar a argumentação.
% \begin{apendicesenv}
% \partapendices
% \chapter{<Título do Apêndice>}
% <Conteúdo do apêndice, como questionários elaborados pela equipe, diagramas extensos, etc.>
% \end{apendicesenv}

% --- ANEXOS ---
% Elemento opcional. Documentos não criados pelo autor que servem de fundamentação ou ilustração.
% \begin{anexosenv}
% \partanexos
% \chapter{<Título do Anexo>}
% <Conteúdo do anexo, como manuais de APIs de terceiros, leis, normas utilizadas, etc.>
% \end{anexosenv}

\end{document}            

% --------------------------------------------------------
% ELEMENTOS PÓS-TEXTUAIS
% --------------------------------------------------------
\postextual

% --- REFERÊNCIAS BIBLIOGRÁFICAS ---
% Elemento obrigatório. Puxa as citações do arquivo referencias.bib formatando nas normas da ABNT (NBR 6023).
\bibliographystyle{abntex2-alf} % Define o estilo ABNT (Autor-Data)
\bibliography{references}

% --- APÊNDICES ---
% Elemento opcional. Documentos criados pelo próprio autor para complementar a argumentação.
% \begin{apendicesenv}
% \partapendices
% \chapter{<Título do Apêndice>}
% <Conteúdo do apêndice, como questionários elaborados pela equipe, diagramas extensos, etc.>
% \end{apendicesenv}

% --- ANEXOS ---
% Elemento opcional. Documentos não criados pelo autor que servem de fundamentação ou ilustração.
% \begin{anexosenv}
% \partanexos
% \chapter{<Título do Anexo>}
% <Conteúdo do anexo, como manuais de APIs de terceiros, leis, normas utilizadas, etc.>
% \end{anexosenv}

\end{document}            

% --------------------------------------------------------
% ELEMENTOS PÓS-TEXTUAIS
% --------------------------------------------------------
\postextual

% --- REFERÊNCIAS BIBLIOGRÁFICAS ---
% Elemento obrigatório. Puxa as citações do arquivo referencias.bib formatando nas normas da ABNT (NBR 6023).
\bibliographystyle{abntex2-alf} % Define o estilo ABNT (Autor-Data)
\bibliography{references}

% --- APÊNDICES ---
% Elemento opcional. Documentos criados pelo próprio autor para complementar a argumentação.
% \begin{apendicesenv}
% \partapendices
% \chapter{<Título do Apêndice>}
% <Conteúdo do apêndice, como questionários elaborados pela equipe, diagramas extensos, etc.>
% \end{apendicesenv}

% --- ANEXOS ---
% Elemento opcional. Documentos não criados pelo autor que servem de fundamentação ou ilustração.
% \begin{anexosenv}
% \partanexos
% \chapter{<Título do Anexo>}
% <Conteúdo do anexo, como manuais de APIs de terceiros, leis, normas utilizadas, etc.>
% \end{anexosenv}

\end{document}            

% --------------------------------------------------------
% ELEMENTOS PÓS-TEXTUAIS
% --------------------------------------------------------
\postextual

% --- REFERÊNCIAS BIBLIOGRÁFICAS ---
% Elemento obrigatório. Puxa as citações do arquivo referencias.bib formatando nas normas da ABNT (NBR 6023).
\bibliographystyle{abntex2-alf} % Define o estilo ABNT (Autor-Data)
\bibliography{references}

% --- APÊNDICES ---
% Elemento opcional. Documentos criados pelo próprio autor para complementar a argumentação.
% \begin{apendicesenv}
% \partapendices
% \chapter{<Título do Apêndice>}
% <Conteúdo do apêndice, como questionários elaborados pela equipe, diagramas extensos, etc.>
% \end{apendicesenv}

% --- ANEXOS ---
% Elemento opcional. Documentos não criados pelo autor que servem de fundamentação ou ilustração.
% \begin{anexosenv}
% \partanexos
% \chapter{<Título do Anexo>}
% <Conteúdo do anexo, como manuais de APIs de terceiros, leis, normas utilizadas, etc.>
% \end{anexosenv}

\end{document}